%%%%% Fichero de ejemplo LaTeX que ilustra el uso de la Hoja de Estilo %%%%%%
%%%%% Jornadas.cls para Jornadas Sarteco

\documentclass[twocolumn,twoside]{Jornadas}
\usepackage[utf8]{inputenc}
\usepackage[spanish]{babel}
\usepackage{listings}
\usepackage{algorithm}
\usepackage{algorithmicx}
\usepackage{algcompatible}
\usepackage{adjustbox}
\usepackage{graphicx}
\usepackage{color}
\usepackage{caption}
\captionsetup{font=footnotesize}
\usepackage[caption=false,font=footnotesize]{subfig}
%\setlength{\marginparwidth}{2cm}
%\usepackage{todonotes}
\usepackage{placeins}
\usepackage{hyperref}
\usepackage{url}

% -*- Ajustes LaTeX en relación a las figuras -*-
\setcounter{topnumber}{10}     % Max. numero de figs. on top
\setcounter{bottomnumber}{10}  % Max. numero de figs. abajo
\setcounter{totalnumber}{10}   % Max. numero de figs. por pagina
\renewcommand{\topfraction}{1} % Max. fraccion de pagina ocupada por figs.
\renewcommand{\bottomfraction}{1}
\renewcommand{\textfraction}{0}  % Min. fraccion de pagina ocupada por texto
\renewcommand{\floatpagefraction}{1} % Max. espacio de pagina solo con figs.

\definecolor{gray97}{gray}{.97}
\definecolor{gray75}{gray}{.75}
\definecolor{gray45}{gray}{.45}

\lstset{
     inputencoding=utf8,
     extendedchars=true,
     backgroundcolor=\color{gray97},
     %
     stringstyle=\ttfamily,
     showstringspaces = false,
     basicstyle=\scriptsize\ttfamily,
     commentstyle=\color{gray45},
     keywordstyle=\bfseries,
     linewidth=.98\columnwidth,
     xleftmargin=3mm,
     breaklines=true,
     numbers=left,
     numbersep=6pt,
     numberstyle=\tiny,
     numberfirstline = false,
     firstnumber=auto,
     breaklines=true,
     %
     escapeinside={(*@}{@*)},
     literate={á}{{\'a}}1 {é}{{\'e}}1 {í}{{\'i}}1 {ó}{{\'o}}1 {ú}{{\'u}}1 {ñ}{{\~n}}1
   }

\def\BibTeX{{\rm B\kern-.05em{\sc i\kern-.025em b}\kern-.08em
    T\kern-.1667em\lower.7ex\hbox{E}\kern-.125emX}}

\newtheorem{theorem}{Teorema}

\hyphenation{pa-ra-le-lis-mo pro-cee-dings}

%Directorios en los que se buscan las figuras
\graphicspath{{.}{./Figuras/}}
%%%%%%%%%%%%%%%%%%%%%%%%%%%%%%%%%%%%%%%%%%%%

\begin{document}

\title{Pipeline de Preprocesamiento y Extracción de Características para Reconocimiento de Patrones}

\author{%
     José Daniel Flores Morales%
     \thanks{Benemérita Universidad Autónoma de Puebla, Facultad de Ciencias de la Computación.}
}

\maketitle
% Oculta las cabeceras y los números de página.
% Ambos elementos se añadirán durante la edición de las actas completas.
\markboth{}{}
\pagestyle{empty} 
\thispagestyle{empty} % Oculta el número de la primera página

\begin{abstract}
Este artículo presenta el diseño y desarrollo de un pipeline integral para el preprocesamiento de datos en tareas de reconocimiento de patrones. El flujo de trabajo abarca la conversión de datos crudos, múltiples técnicas de normalización (Min-Max, Z-Score, Robust Scaler y Decimal Scaling), y la extracción de características mediante la correlación de Pearson, Análisis de Componentes Principales (PCA) y Análisis de Correlación Canónica (CCA). Finalmente, se implementó una estrategia de partición estratificada de los datos (Train, Validation, Test) para garantizar la robustez en la fase posterior de entrenamiento de modelos predictivos.
\end{abstract}

\begin{keywords}
Reconocimiento de Patrones, Normalización, PCA, CCA, Correlación de Pearson, Preprocesamiento.
\end{keywords}

%Añade los ficheros que necesites:

\section{Introducción}
El reconocimiento de patrones y el aprendizaje automático dependen intrínsecamente de la calidad de los datos de entrada \cite{bishop2006,fieguth2022}. En muchos casos empíricos, los datos crudos presentan ruido, escalas dispares y alta dimensionalidad que obstaculizan el rendimiento de los modelos matemáticos y probabilísticos \cite{duda2000}.

Por consiguiente, el preprocesamiento de datos se erige como una fase crítica. Este artículo documenta la construcción de un pipeline estructurado en cuatro etapas principales: procesamiento inicial, normalización, análisis de correlación (reducción de dimensionalidad) y partición de datos. El objetivo principal es transformar un conjunto de datos en bruto en subconjuntos de entrenamiento, validación y prueba óptimos, garantizando la eliminación de redundancias y la escalabilidad del modelo.

El presente estudio evalúa empíricamente diversas estrategias matemáticas sobre dos conjuntos de datos multidimensionales (\textit{Slice} y \textit{Vh}), detallando la toma de decisiones algorítmicas en cada etapa del desarrollo.

\section{Marco Teórico}
El desarrollo del pipeline fundamenta su arquitectura en técnicas establecidas de la ciencia de datos y el reconocimiento estadístico de patrones \cite{jain2000}.

\subsection{Normalización de Datos}
La normalización proyecta los datos a un espacio uniforme, evitando que variables con escalas mayores dominen el cálculo de distancias (ej. en clasificadores KNN). Se implementaron los siguientes métodos:
\begin{itemize}
    \item \textbf{Min-Max:} Ajusta los datos en un rango de [0,1].
    \item \textbf{Z-Score (Estandarización):} Centra los datos en una media de 0 y desviación estándar de 1, útil para algoritmos que asumen distribución normal.
    \item \textbf{Robust Scaler:} Emplea la mediana y el rango intercuartílico, mitigando el impacto de valores atípicos (outliers).
    \item \textbf{Decimal Scaling:} Escala los datos desplazando el punto decimal según el valor absoluto máximo.
\end{itemize}

\subsection{Reducción de Dimensionalidad}
La \textit{Maldición de la Dimensionalidad} provoca que los clasificadores pierdan eficiencia ante un exceso de características. Para mitigarla, se utilizaron tres enfoques:
\begin{itemize}
    \item \textbf{Correlación de Pearson:} Mide la dependencia lineal perfecta entre dos variables. Se utiliza para eliminar características redundantes (clones) que comparten un alto coeficiente ($|r| > 0.98$).
    \item \textbf{Análisis de Componentes Principales (PCA):} Transformación ortogonal que proyecta los datos hacia componentes que maximizan la varianza.
    \item \textbf{Análisis de Correlación Canónica (CCA):} Busca relaciones lineales que maximizan la correlación entre dos conjuntos de variables (en este caso, características contra etiquetas de clase).
\end{itemize}

\section{Metodología}
La arquitectura del proyecto se diseñó modularmente mediante un orquestador \texttt{main.py} que ejecuta secuencialmente las cuatro etapas.

\subsection{Etapa 1: Preprocesamiento}
Los conjuntos de datos \textit{Slice.txt} y \textit{Vh.txt} fueron convertidos a un formato estandarizado \texttt{.csv} (\textit{Comma Separated Values}) para facilitar la carga vectorial mediante la biblioteca Pandas.

\subsection{Etapa 2: Normalización}
La normalización se aplicó a las matrices de características de ambos datasets, generando tres versiones transformadas por archivo: \textit{Min-Max}, \textit{Z-Score} y \textit{Robust}. Posteriormente, se seleccionó empíricamente la versión Z-Score como línea base para la reducción de dimensionalidad por su eficacia en evitar sesgos de escala.

\subsection{Etapa 3: Análisis y Reducción}
A partir del dataset estandarizado con Z-Score, se ejecutaron tres procesos paralelos de reducción:
\begin{enumerate}
    \item \textbf{Selección de Características:} Se calculó la matriz de covarianza usando la correlación de Pearson. Recorriendo la matriz triangular superior, se identificaron y eliminaron aquellas características con una correlación $|r| > 0.98$.
    \item \textbf{Proyección PCA:} Se aplicó PCA para reducir el conjunto a 10 componentes principales, registrando el porcentaje de varianza conservada.
    \item \textbf{Transformación CCA:} Se aplicó CCA para maximizar la correlación entre las características continuas y las etiquetas discretas.
\end{enumerate}

\subsection{Etapa 4: División de Datos}
Para prevenir la filtración de datos (\textit{Data Leakage}) durante el futuro entrenamiento, los datos reducidos se partieron utilizando un muestreo estratificado (\texttt{stratify=y}) para preservar la distribución de clases original. La división fue: $70\%$ Entrenamiento (\textit{Train}), $15\%$ Validación (\textit{Validation}) y $15\%$ Prueba (\textit{Test}).
\section{Resultados y Experimentación}
La ejecución del pipeline demostró reducciones sustanciales en la complejidad de los conjuntos de datos.

\subsection{Impacto del Análisis de Correlación}
Al someter los datasets al filtro de correlación de Pearson ($|r| > 0.98$):
\begin{itemize}
    \item \textbf{Dataset Slice:} Se detectaron 148 pares altamente correlacionados. Al descartar la redundancia, el volumen original de 385 columnas se redujo a 324 características limpias (eliminación de 61 "clones").
    \item \textbf{Dataset Vh:} De 18 columnas originales, se detectaron 3 pares correlacionados, generando un dataset limpio final de 16 columnas.
\end{itemize}

\subsection{Rendimiento de PCA}
La reducción explícita a 10 componentes mediante PCA arrojó los siguientes niveles de varianza retenida:
\begin{itemize}
    \item En \textit{Slice}, las 10 componentes lograron conservar el $71.99\%$ de la información original.
    \item En \textit{Vh}, las 10 componentes consiguieron encapsular el $98.90\%$ de la varianza total, evidenciando una alta redundancia intrínseca en el espacio original.
\end{itemize}

\subsection{Distribución de Datasets Resultantes}
Finalmente, la separación en conjuntos estratificados generó archivos listos para el entrenamiento. Para el caso del dataset \textit{Slice\_zscore\_pearson\_limpio}:
\begin{itemize}
    \item \textbf{Train (70\%):} 286 muestras.
    \item \textbf{Validation (15\%):} 61 muestras.
    \item \textbf{Test (15\%):} 62 muestras.
\end{itemize}
La misma distribución se garantizó para todas las variantes (\textit{pearson, pca, cca}) de los datasets procesados.
\section{Conclusiones}
La implementación de este pipeline estructurado ha logrado sistematizar el preprocesamiento de grandes volúmenes de datos numéricos. Se comprobó que el análisis de correlación de Pearson es altamente efectivo para la selección de características lineales puras, mientras que métodos como PCA permiten una compresión agresiva del espacio dimensional conservando una alta densidad de información (hasta $98.90\%$ en el dataset Vh).

Contar con datasets rigurosamente normalizados, purgados de variables redundantes y estratificados metodológicamente garantiza que los modelos de clasificación posteriores operen con la máxima eficiencia espacial y computacional posible, previniendo sobreajustes y asegurando la validez del modelo final.


\section*{Agradecimientos}

Añadir en esta sección no numerada los posibles agradecimientos a proyectos,
instituciones o personas. Por ejemplo: El presente trabajo ha sido financiado
mediante  el proyecto xxxx-xxxxx-xxx. 

%Comenta la línea \nocite para que sólo se incluyan en las referencias
%las entradas que esté referenciadas en el texto con \cite{}
\nocite{*}
\bibliographystyle{Jornadas}
\bibliography{biblio}

\end{document}

